% Converted from the authoritative recovery source, 2026-11-14.
\subsection{Extension to finitely many exchangeable
species}\label{influence-recovery:extension-to-finitely-many-exchangeable-species}

Let \([n]=B_1\sqcup\cdots\sqcup B_r\) be a known partition, and suppose
\(\mu\) is invariant under separate permutations inside each species.
Conditional on the vector of species weights, it is a product of uniform
slices; the weights themselves may be arbitrarily dependent. There are
absolute constants such that the same potential satisfies \[
 \mathcal J_\mu(f;\mathcal T)\ge c v^{7/2}/d^7
 \quad\text{with depth }O(r^3d^{18}/v^7).                 \tag{20}
\] Consequently the influence iteration gives
\(O(r^3d^{27}\epsilon^{-21})\) deterministic queries. For fixed \(r\),
this is dimension independent and uniform over all species sizes and
mixing laws. If \(r\) grows with \(n\), the displayed dependence on
\(r\) must be retained.

First extend (6) to arbitrary coordinate-bias vectors. For a globally
bounded degree-\(s\) polynomial \(h\), write
\(F(\mathbf p)=\mathbb E_{\pi_{\mathbf p}}h\). Differentiating this
multilinear product expectation gives \[
 \sum_i|\partial_{p_i}F(\mathbf p)|
 \le\mathbb E_{\pi_{\mathbf p}}
      \sum_i|h(X^{\oplus i})-h(X)|\le4s^2.               \tag{21}
\] The last bound is pointwise: use the signed-sum Markov argument
separately for positive and negative differences of \(h\). This also
holds at boundary biases by polynomial continuity. Integrating along the
segment from \(\mathbf p\) to \(\mathbf q\) yields \[
 |\mathbb E_{\pi_{\mathbf p}}h-
       \mathbb E_{\pi_{\mathbf q}}h|
 \le4s^2\|\mathbf p-\mathbf q\|_\infty.                 \tag{22}
\] Applying (22) to \(f,f^2,(D_i f)^2\) bounds the corresponding
variance and influence errors by respectively \(24d^2\eta\) and
\(16d^2\eta\), where \(\eta=\|\mathbf p-\mathbf q\|_\infty\). Round
every bias to a dyadic with the same denominator \(2^b\). The source
comparison (8) still holds, with a separate threshold gadget for each
coordinate. Increasing the absolute constant \(B\) in the choice of
\(b\) absorbs the changed constants in (9). Thus (6), with adjusted
absolute constants, holds for every product law \(\pi_{\mathbf p}\),
including arbitrary endpoint biases.

For the species theorem, choose a prefix length \(L\). Read in full
every species of size less than \(2L\); in each other species read a
fixed prefix of \(L\) coordinates. This uses at most \(2rL\) queries.
Let \(H\) be the whole transcript. In each unread species \(b\), write
\(N_b\ge L\) for its remaining size, \(q_b\) for its observed prefix
density, and \(P_b\) for its latent remaining density. Define \[
 b_H=\sum_{b\text{ unread}}
       \left(\mathbb E[|P_b-q_b|\mid H]+\frac1{2\sqrt{N_b}}\right).
\] Conditional on the original vector of species weights, each prefix
count is hypergeometric, so (12) holds in each species. Averaging gives
\(\mathbb E_H b_H\le3r/(2\sqrt L)\), regardless of dependence among the
weights.

For a fixed remaining weight vector, telescope the replacement of the
individual slice factors by their bias-\(q_b\) products. At each step,
integrate the polynomial over all other species; the resulting function
on the current species remains globally bounded and has degree at most
the original degree. Thus (11) gives, after conditioning on \(H\) and
averaging the weight vector, \[
 |\mathbb E_{\mu_H}h-\mathbb E_{\pi_{\mathbf q}}h|
 \le4s^2b_H.                                           \tag{23}
\] This argument permits correlated latent weights: the telescoping is
performed for each weight vector before it is averaged. Consequently
(14) holds with this new \(b_H\), and the average error is at most
\(e_L=36rd^2/\sqrt L\).

Now use the same \(h_v\) as in (16), with the absolute product constant
just established, and choose \(L=\lceil(72rd^2/h_v)^2\rceil\). Then
\(e_L\le h_v/2\). If no unread coordinates remain, the transcript
determines \(f\). Otherwise repeat (17)--(18): either the initial
transcript reveals variance at least \(v/2\), or a coordinate selected
from it has expected conditional derivative influence at least
\(h_v/2\). Query that coordinate. The depth is at most \(2rL+1\),
proving (20). Pinnings preserve the species symmetries, so the preceding
potential iteration applies unchanged, with only the factor \(r^3\) in
its per-round query cost.
